\documentclass[11pt]{article}
\usepackage[utf8]{inputenc}
\usepackage[T1]{fontenc}
\usepackage[margin=2.4cm]{geometry}
\usepackage{amsmath,booktabs,siunitx}
\usepackage{enumitem}
\usepackage{fancyhdr}
\usepackage[hidelinks]{hyperref}
\pagestyle{fancy}\fancyhf{}\lhead{Reporte de laboratorio}\rhead{\thepage}
\setlist{nosep,leftmargin=*}
\title{Reporte de Laboratorio}\author{Nombre \hspace{2cm} Matrícula}
\date{\today}
\begin{document}
\maketitle
\begin{abstract}
Resuma aquí el propósito, el método, el resultado principal y la conclusión
en no más de 150 palabras.
\end{abstract}
\section{Introducción}
Presente el fenómeno, la pregunta de investigación y la hipótesis. Defina las
variables independientes, dependientes y de control.
\section{Fundamento y método}
Explique el modelo utilizado, por ejemplo
\[
  \bar{x}=\frac{1}{n}\sum_{i=1}^{n}x_i,\qquad
  s=\sqrt{\frac{1}{n-1}\sum_{i=1}^{n}(x_i-\bar{x})^2}.
\]
\begin{enumerate}
  \item Liste el equipo y sus resoluciones.
  \item Describa el procedimiento en orden suficiente para reproducirlo.
  \item Registre condiciones ambientales y observaciones inesperadas.
\end{enumerate}
\section{Resultados}
\begin{table}[ht]
\centering
\caption{Tabla de ejemplo para sustituir por datos experimentales.}
\begin{tabular}{@{}lrr@{}}\toprule
Muestra & Medición & Incertidumbre \\ \midrule
A & 12.4 & 0.3 \\
B & 15.1 & 0.4 \\
C & 14.8 & 0.3 \\ \bottomrule
\end{tabular}
\end{table}
\section{Discusión}
Compare las mediciones con la predicción, cuantifique el error relativo e
indique qué observación respalda o contradice la hipótesis.
\section{Conclusión}
Responda la pregunta inicial y proponga una mejora concreta para el siguiente
experimento.
\end{document}
